\chapter{图形学与交互技术}

\section{领域导言与演进背景}

\begin{preamblebox}[知识领域概述与变革]
计算机图形学是利用计算机生成与处理数字图像的核心学科，其本质可以被视为“以计算为媒介实现视觉感知与人机信息沟通的科学”。从最初的线框显示到当今照片级拟真画质，计算机图形学的应用已渗透至现代计算的方方面面：机器学习、医学成像、工程辅助设计、科学与信息数据可视化、三维动画电影、电影视觉特效（VFX）、沉浸式训练仿真器以及高交互数字游戏等。

在传统的计算机科学本科教育中，计算机图形学主要聚焦于三维渲染管线、线性代数与空间变换、几何光学物理现象学模拟等经典主题。然而在现代高级培养阶段，全球顶尖高校越来越倾向于开设多门高度分化且面向前沿的进阶专向课程，例如游戏工程引擎架构、高级角色动画、科学可视化、物理/实体计算以及涵盖虚拟现实、增强现实与混合现实的沉浸式计算。该领域与其他计算机专业核心领域存在极其深厚的交叉与重叠，尤其是：人工智能（AI）、人机交互（HCI）、并行与分布式计算（PDC）、专用平台开发（SPD）以及社会、伦理与专业素养（SEP）。

为了使计算机科学专业学生能够娴熟地运用并开发计算机图形系统，教学必须深入剖析一系列具体的工程与底层实现机制，包括人类视觉感知与认知心理学、数字图像与几何文件编码规范、底层硬件抽象接口以及现代图形渲染应用程序接口。通过理论推导与工程实践的紧密结合，学生将具备作为高阶用户和底层开发者来理解、评估、重构并实现前沿图形与交互技术的核心能力。鉴于该领域技术演进一日千里，CS2023 工作组在制定本大纲时特意避免过度限定具体商业 API 或特定编程语言，文中所提及的开发接口、图形系统与开发工具代表了截至 2022 年前后的典型工业界与学术界主流范式。

\textbf{自 CS2013 以来的核心演进与重大变革：}
\begin{enumerate}
    \item \textbf{领域命名与 ACM SIGGRAPH 国际标准对齐}：为更契合现代前沿计算形态，CS2023 将 CC2013 中的“图形学与可视化 (GV)”正式更名为“图形学与交互技术 (GIT)”，实现与计算机图形学全球顶级学术组织 ACM SIGGRAPH 的命名统一。
    \item \textbf{移动端普及与算力下沉驱动的全民视觉化}：自 2013 年以来，机器学习、计算机视觉、数据科学与 AI 的爆发，以及从智能汽车到家用物联网设备广泛内嵌的传感器，深刻改变了图形交互的生态。如今移动智能手机的算力飞跃，使全球数十亿普通大众既是高画质图形与交互的消费者，更是创作者。2013 年仅能在高性能桌面图形工作站上运行的复杂动画骨骼蒙皮、流体物理仿真、实时全局光照与沉浸式应用，如今已能在手持设备上流畅运行。
    \item \textbf{泛在图像与文化产业数字化重塑}：十年来全球数字数据爆发式增长，新闻媒体与科研机构每天都在利用高维信息可视化揭示自然与社会规律；移动数字游戏的全球产值已经超越电影与音乐两大传统娱乐工业的总和；而在现代电影工业中，计算机生成图像（CGI）已成为几乎所有商业影视的不可或缺之基石。
    \item \textbf{全面扩充五大前沿核心知识单元}：为捕捉学科外延的剧烈拓展，在原有知识单元（基本概念、基础渲染、几何建模、高级着色、计算机动画、数据可视化）之上，全新增设了以下 5 个知识单元：
    \begin{itemize}
        \item \textbf{沉浸式技术 (Immersion: MR, AR, VR)}：系统建立空间计算、动光延迟抑制与虚实融合机制。
        \item \textbf{交互技术}：涵盖事件驱动范式、多模态融合、力反馈外骨骼与脑机接口前沿。
        \item \textbf{数字图像处理}：融合形态学、空频域滤波、图像修复以及深度卷积与 Transformer 视觉前沿。
        \item \textbf{实体与物理计算}：贯通传感器采集、执行器控制、数控数字制造（CAD/CAM/3D打印）与物联网系统。
        \item \textbf{计算机物理仿真}：确立面向图形视觉可信度的质点弹簧、SPH 粒子流体、欧拉网格烟火与碰撞动力学。
    \end{itemize}
    \item \textbf{深度审视图形与视觉技术的伦理困境}：人工智能生成内容（AIGC）与高保真图形合成技术的结合，带来了前所未有的伦理考验。师生必须直面在自动驾驶机器视觉感知失效、深空遥感伪影、计算摄影学“脑补”失真、基于人脸与身形的深度伪造（Deepfakes）以及人脸识别在不同种族发型上的算法偏见与不公等现实伦理痛点。
\end{enumerate}
\end{preamblebox}

\section{核心学时分配建议}

在 CS2023“向日葵模型”架构下，图形学与交互技术领域各单元的核心推荐学时分配如下表所示：

\begin{table}[H]
\centering
\setlength{\tabcolsep}{3.5pt}\fontsize{8pt}{9.6pt}\selectfont\begin{tabularx}{\columnwidth}{Y c c}
\toprule
\textbf{知识单元} & \textbf{计算机核心} & \textbf{领域核心} \\
\midrule
基本概念 & 4 & -- \\
基础渲染技术 & -- & 18 \\
几何建模 & -- & 6 \\
高级着色与进阶渲染 & -- & 6 \\
计算机动画 & -- & 6 \\
数据与信息可视化 & -- & 6 \\
沉浸式技术 & -- & 6 \\
交互技术 & -- & 6 \\
数字图像处理 & -- & 6 \\
实体与物理计算 & -- & 6 \\
计算机物理仿真 & -- & 6 \\
\midrule
\textbf{合计} & \textbf{4 学时} & \textbf{72 学时} \\
\bottomrule
\end{tabularx}
\end{table}

\section{知识单元与详细大纲}

\subsection{计算机图形学基本概念}
\textbf{【学时建议】}：4 计算机核心学时。

对于几乎所有计算机专业毕业生与软件开发者而言，理解人机视觉交互的底层规律都是必不可少的素养。本单元的核心内容既可作为独立图形学课程的先导篇章，亦可融入大学计算机导论或软件开发基础（SDF）课程中讲授。

\subsubsection{计算机核心主题}
\begin{itemize}
    \item \textbf{娱乐、工商业与科学应用图景}：视觉特效（VFX）、机器学习与计算机视觉交叉、现代图形用户界面（GUI）、视频编辑管线、游戏引擎架构、计算机辅助设计与制造（CAD/CAM）、空间与时变数据可视化、虚拟/增强/混合现实（VR/AR/MR）。
    \item \textbf{人类视觉生理系统与感知模型}：
    \begin{itemize}
        \item 视网膜感光锥状细胞与三刺激响应（RGB 三原色理论）。
        \item 经典“眼球作为摄像机”的透视投影成像几何。
        \item 视觉暂留特性、临界闪烁融合频率、帧率与动态模糊。
        \item 对比度感知机制：边缘增强特性、马赫带效应、色彩抖动与空间走样。
        \item 人眼对光强的非线性对数响应：高动态范围（HDR）与色调映射必要性。
        \item 双目立体视觉（Binocular Vision 与双目视差 Stereo）。
        \item 视觉无障碍包容性设计：红绿色盲等色觉辨认障碍、频闪光敏性癫痫规避、单目弱视交互补偿（与 SEP、HCI 深度交叉）。
    \end{itemize}
    \item \textbf{模拟物理信号的数字化机制}：
    \begin{itemize}
        \item 光栅图元表征与矢量数学表征的本质对立与时空权衡：多边形面片、体积标量场与连续实体对象。
        \item 分辨率度量：显示屏幕的像素密度（PPI）与激光打印设备的点密度（DPI）。
        \item 连续空间的离散空间采样与色彩强度量化。
    \end{itemize}
    \item \textbf{标准媒体编码与格式}：
    \begin{itemize}
        \item 光栅图像格式：无损压缩编码（如 PNG、无损 TIFF）与有损压缩编码（如 JPEG 离散余弦变换）。
        \item 矢量数据格式规范（如 SVG、PDF 路径）。
    \end{itemize}
    \item \textbf{经典色彩模型与空间}：加色模型（RGB，用于自发光显示器）、减色模型（CMYK，用于印刷油墨染料）、面向人类感知直觉的色彩模型（HSV/HSL，色相、饱和度、明度）。
    \item \textbf{存储空间占用与动态重算时延的权衡}：光栅高保真缓存与矢量数学解析重绘的计算复杂度对比。
    \item \textbf{动画的离散本质}：连续时空运动在离散时序静止画面帧序列中的表征机理。
    \item \textbf{伦理、版权与社会议题 (SEP)}：换脸与深度伪造（Deepfakes）技术、人脸识别中的偏见、个人生物特征隐私合规与视觉作品知识产权保护。
\end{itemize}

\subsubsection{示范性学习成果与能力目标}
\begin{enumerate}
    \item \textbf{[识别]} 列举并系统剖析向人类用户呈现数字信息的通用技术方案（包括计算机图形学与声学技术）。
    \item \textbf{[解释]} 阐述连续模拟物理信号如何通过离散采样合理表征为数字信号（例如连续光学图像如何量化表征为二维像素阵列）。
    \item \textbf{[计算]} 给定目标分辨率、颜色深度与通道数，精确计算未压缩全彩光栅图像所占用的物理内存容量。
    \item \textbf{[开发]} 绘制或编程构建图表，演示人类视觉感知系统的生理极限如何直接制约模拟信号数字化采样参数的工程决策。
    \item \textbf{[设计]} 针对标准人机交互界面，设计一套能够适应红绿色盲或色弱用户的替代色彩与形状双编码无障碍界面方案。
    \item \textbf{[开发]} 调用行业标准图形 API（如 HTML5 Canvas、Qt 或基础 OpenGL），构建一个包含交互图元的简易图形用户界面。
    \item \textbf{[权衡]} 论证在何种工程情境下应当选用 JPG、PNG、MP3、MP4 或 GIF 文件格式，并从压缩机理与感知失真角度给出深度技术论据。
    \item \textbf{[评估]} 指出通用媒体格式中采用的典型有损与无损压缩算法实例，并分析其香农信息论与心理物理学依据。
    \item \textbf{[解释]} 系统阐释常见色彩空间（RGB、CMYK、HSV）的坐标基底，并说明其在物理显示输出设备中的硬件映射关系。
    \item \textbf{[分析]} 深入辨析在图像存储中“完整持久化光栅数据”与“仅存储关键几何信息并在渲染时动态重建”之间的时空权衡边界。
    \item \textbf{[计算]} 给定时长 $n$ 秒、帧率 $f$ 帧/秒及色彩分辨率 $r$，计算未压缩视频流的实时内存带宽与物理存储开销。
    \item \textbf{[解释]} 阐释人类视觉生理机制中从离散静态定格画面序列生成连续平滑运动感知的基本过程（临界闪烁融合机理）。
    \item \textbf{[解释]} 举例说明对现实模拟物理世界进行数字离散采样时可能出现的视觉虚假走样效应（如莫尔条纹、车轮倒转效应）。
\end{enumerate}

\subsection{数据与信息可视化}
\textbf{【学时建议】}：6 领域核心学时。

可视化是为抽象数据建立视觉图形表征的核心过程。本单元与人机交互（HCI）及计算科学（CN）存在极强的内在关联，涉及目标受众的认知特征评估与交互界面的效能测度。

\subsubsection{领域核心主题}
\begin{itemize}
    \item \textbf{数据可视化与信息可视化辨析}：科学计算数据场（强空间绑定）与抽象结构化信息（非物理空间实体）的可视化差异。
    \item \textbf{多模态数据类型的可视化表征}：
    \begin{itemize}
        \item 二维/三维标量场数据：等值线、等值面、伪彩色映射。
        \item 矢量场与流场数据：方向线、流线、脉线、迹线与粒子轨迹跟踪。
        \item 时变时序多维动态数据场。
        \item 超高维多元特征数据：散点图矩阵（SPLOM）、平行坐标轴。
        \item 非空间关系型数据：层次树状图、力导向网络拓扑图、词云与文本语义演化流。
    \end{itemize}
    \item \textbf{高阶可视化变换技术}：高维降维投影（PCA 主成分分析、t-SNE、UMAP）、多元数据空间剖分与交互联动刷选（Brushing \& Linking）。
    \item \textbf{直接体绘制算法}：光线投射算法、传输函数的色彩与不透明度分配、体素分割与边缘增强。
    \item \textbf{通用科学与空间数据存储规范}：HDF5、netCDF、GeoTIFF、原始二进制连续流、结构化 CSV、规范化文本序列。
    \item \textbf{工业界与学术界主流可视化工具链}：R 语言（ggplot2）、Processing、D3.js、地理信息系统（GIS）、MATLAB、IDL、Python 可视化生态（Matplotlib、Seaborn、VTK）。
    \item \textbf{驱动视觉抽象的认知心理学底座}：视觉感知信道理论（位置、长度、角度、面积、颜色饱和度）、贝特定律、格式塔心理学原则、色彩感知调色板（感知均匀色阶如 Viridis、区分性色阶）。
    \item \textbf{端到端可视化工程设计原则}：
    \begin{itemize}
        \item 目标导向性：科研探索与公众科学传播的区别。
        \item 受众分层定位：专业领域科学家与大众公众。
        \item 负责任与道德规范的可视化：防范虚假视觉修辞（坐标轴刻度截断欺骗、过度平滑消除关键极值、数据清洗阶段的恶意填补插值）；识别真数据引起的视觉误导（如高频运动引起的频闪混叠导致风扇叶片视觉停转）。
    \end{itemize}
    \item \textbf{可视化有效性评估方法}：任务完成时间、准确率眼动追踪分析与主观认知负荷评估量表。
    \item \textbf{可视化认知偏差与偏见}：知觉锚定效应、框架效应与视觉突显偏差。
\end{itemize}

\subsubsection{示范性学习成果与能力目标}
\begin{enumerate}
    \item \textbf{[开发]} 编程实现经典基础可视化算法（如基于移动立方体 Marching Cubes 的等值面提取、离散标量场传递函数着色映射）。
    \item \textbf{[评估]} 在计算时间开销、显存带宽占用与视觉呈现保真度之间，量化评估不同可视化算法的工程权衡。
    \item \textbf{[设计]} 针对特定领域的数据拓扑特征、具体分析任务与目标受众类型，制定兼具美学价值与认知严谨性的可视化设计系统方案。
    \item \textbf{[分析]} 剖析给定工业级或科学可视化图表在执行特定决策任务时的有效性与认知效率，指出潜在设计缺陷。
    \item \textbf{[评估]} 设计一套严谨的用户体验实验流程，以定量与定性手段科学评测某可视化算法或软件系统的实际效用。
    \item \textbf{[识别]} 调研并归纳可视化技术在现代医学断层成像诊断、气象流场预警、天体物理模拟及金融风险空间地理分析中的应用范式。
\end{enumerate}

\subsection{基础渲染技术}
\textbf{【学时建议】}：18 领域核心学时。

本单元涵盖计算机图形学的核心物理与数学基础，是所有高阶计算机图形学、游戏开发与三维工程软件的共同理论基石。

\subsubsection{领域核心主题}
\begin{itemize}
    \item \textbf{现代计算机图形渲染管线}：输入装配、顶点着色、几何与曲面细分处理、图元剪裁与屏幕映射、光栅化、片元/像素着色、深度测试与模板测试、Alpha 混合及帧缓冲输出。
    \item \textbf{自然界光传输物理学与数值积分}：辐射通量、辐照度与辐亮度；光源辐射、介质散射、能量守恒定律；渲染方程的半球立体角积分形式及其向离散数值积分的转换。
    \item \textbf{正向渲染与反向渲染架构}：基于物体空间投影的光栅化与基于图像空间像素追踪的光线投射之范式对比与计算复杂度对偶性。
    \item \textbf{三维多边形网格表征}：三角网格拓扑、顶点缓冲区（VBO）、顶点索引缓冲区（EBO）、顶点法向量与切线空间。
    \item \textbf{几何光学透视与针孔投影模型}：相似三角形原理、透视除法、视场角（FOV）与平截头体构建。
    \item \textbf{仿射变换与三维齐次坐标系系统}：
    \begin{itemize}
        \item 局部坐标系、世界坐标系、观察/相机坐标系、裁剪坐标系与规范化设备坐标系（NDC）。
        \item 平移、旋转（Rodrigues 公式）、缩放矩阵以及 $4\times 4$ 齐次变换矩阵的非交换性链式复合。
    \end{itemize}
    \item \textbf{光线追踪基础算法}：主光线生成、光线与几何图元（球体、平面、三角形）求交、阴影感知射线检测、理想镜面反射与折射（Snell 定律与 Fresnel 方程）。
    \item \textbf{可见性计算与遮挡消隐机制}：深度缓冲算法（$Z$-buffer / Depth Buffering）、画家算法（Painter's Algorithm 深度排序局限性）、基于光线穿透的消隐判断。
    \item \textbf{三角形基元光栅化算法}：基于平面半空间测试与重心坐标的图元覆盖率判定与属性线性透视校正插值。
    \item \textbf{基于现代可编程着色器 API 的工程渲染}：现代可编程管线架构（如 OpenGL 3.3+ Core Profile、Vulkan 或 DirectX），着色语言编程与Uniform缓冲对象。
    \item \textbf{纹理映射与采样滤波理论}：UV 纹理坐标映射、纹理包裹模式、双线性插值与双三次插值、纹理放大与缩小走样、三线性多级渐远纹理映射与各向异性过滤。
    \item \textbf{空间加速数据结构在渲染中的应用}：轴对齐包围盒（AABB）、包围盒层次结构（BVH）、八叉树与 KD-Tree 在光线求交与视锥剔除中的应用。
    \item \textbf{采样理论与反走样技术}：香农采样定理、奈奎斯特采样频率极限、超采样反走样（SSAA）、多重采样反走样（MSAA）与时间序列抗锯齿（TAA）初步。
    \item \textbf{层次化场景图模型}：树状拓扑层次关系、局部变换向全局变换的递归遍历累积。
\end{itemize}

\subsubsection{示范性学习成果与能力目标}
\begin{enumerate}
    \item \textbf{[解释]} 绘制光线能量在物理场景中从光源发射、表面复杂多重散射并最终被虚拟镜头接收的传输图解，阐明其与半球数值积分的数理对应关系。
    \item \textbf{[描述]} 详细叙述现代硬件图形渲染管线的各个核心阶段及其在 GPU 硬件流水线中的协作逻辑。
    \item \textbf{[对比]} 深入对比正向渲染（光栅化管线）与反向渲染（光线追踪）在处理海量动态光源与复杂几何遮挡时的时空算力开销差异。
    \item \textbf{[开发]} 编写独立程序，在图形窗口中高效绘制二维动态矢量几何形状。
    \item \textbf{[开发]} 基于可编程渲染管线，编写并运行能够加载并显示交互式三维网格模型的渲染系统。
    \item \textbf{[推导]} 利用相似三角形几何原理，推导出将三维坐标点 $(x, y, z)$ 投射映射到投影平面点 $(x/z, y/z, 1)$ 的线性透视除法数学公式。
    \item \textbf{[计算]} 给定由平移、旋转与缩放组成的复合仿射矩阵，精确计算二维与三维几何顶点经矩阵变换后的新坐标。
    \item \textbf{[应用]} 掌握三维笛卡尔坐标系统，推演将二维变换算法扩展至三维齐次变换矩阵所需的维数扩展与基底变换法则。
    \item \textbf{[解释]} 深刻阐释纹理映射的基本原理，以及欠采样与过采样诱发走样的物理机理和 MIP-mapping 的消除策略。
    \item \textbf{[评估]} 在处理复杂三维几何遮挡与可见性判断问题时，系统横向评测光线追踪与 $Z$-buffer 算法的优劣。
    \item \textbf{[开发]} 编程实现对二维几何图元的平移、旋转、缩放复合变换以及齐次视锥剪裁（如 Cohen-Sutherland 或 Liang-Barsky 算法）。
    \item \textbf{[开发]} 基于现代着色器接口（如现代 OpenGL / WebGL），利用顶点缓冲（VBO）、索引缓冲（EBO）与编写 GLSL 顶点/片元着色器，实现一个完整的实时着色器渲染器。
    \item \textbf{[对比]} 深入对比 Phong 经验光照模型、Gouraud 顶点光照插值与基于物理的 Microfacet 材质渲染在真实感上的阶跃演进。
    \item \textbf{[计算]} 根据屏幕显示分辨率、每像素颜色位深与通道格式，精确计算双缓冲图形输出所需的物理显存带宽与容量。
    \item \textbf{[计算]} 根据显示器垂直刷新率（如 60Hz、144Hz）与几何三角形吞吐量，量化推导光栅化渲染管线在单帧时间预算（如 16.6ms、6.9ms）内的算力瓶颈。
\end{enumerate}

\subsection{几何建模}
\textbf{【学时建议】}：6 领域核心学时。

几何建模涵盖二维形状与三维立体几何实体的数学表征、程序化创建、拓扑交互与编辑操作。

\subsubsection{领域核心主题}
\begin{itemize}
    \item \textbf{空间几何原语计算与邻近性测试}：点-线-面三维求交、射线求交、距离场计算、包围盒快速相交与分离轴判定定理（SAT）。
    \item \textbf{三维表面表征与网格几何模型}：
    \begin{itemize}
        \item 几何细分与镶嵌（Tessellation，曲面多边形化）。
        \item 网格拓扑数据结构（半边结构 Half-Edge、翼边结构 Winged-Edge）、网格光顺与平滑算法。
        \item 经典网格生成技术：平面点集 Delaunay 三角剖分、三维隐式等值面提取移动立方体算法。
        \item 参数化多项式曲线与曲面：Hermite 曲线、Bézier 曲线与曲面特征多边形网格、B-样条与非均匀有理 B-样条（NURBS）。
        \item 隐式几何曲面（Implicit Surfaces，代数曲面方程、融合电位球 Metaballs）。
        \item 空间二叉分割树（BSP 树）与场景几何划分。
    \end{itemize}
    \item \textbf{体素与体积几何模型}：
    \begin{itemize}
        \item 三维规则体素网格与三维点云数据模型。
        \item 有符号距离场 (SDF)解析与网格交互融合。
        \item 层次化稀疏体素网格架构（如 OpenVDB 拓扑与存储优化）。
        \item 构造实体几何 (CSG)：正方体、球体、圆柱体等基本几何实体的布尔交、并、差集合运算。
    \end{itemize}
    \item \textbf{过程式与生成式几何建模}：
    \begin{itemize}
        \item 分形几何（Fractals，自相似性 Mandelbrot 集合、分形地形生成）。
        \item 文法驱动的生成式系统：L-系统（Lindenmayer Systems，植物分支形态生成，与编程语言形式文法交叉）。
        \item 基于深度学习的生成式几何模型（如 NeRF 神经辐射场与 3D 高斯泼溅 Gaussian Splatting 初步）。
    \end{itemize}
    \item \textbf{弹性变形与自由形态变形}：
    \begin{itemize}
        \item 物理准静态几何变形方法。
        \item 双调和曲面形变场捕获。
    \end{itemize}
    \item \textbf{多分辨率网格模型与几何简化}：多细节层次（LOD）、基于二次误差度量（QEM）的边折叠网格简化算法。
    \item \textbf{三维反求几何重建技术 (3D Reconstruction)}：基于激光雷达点云与多视角立体视觉（MVS）的表面拓扑重构算法。
\end{itemize}

\subsubsection{示范性学习成果与能力目标}
\begin{enumerate}
    \item \textbf{[对比]} 深入对比参数化曲面（如 Bézier / NURBS）与隐式代数曲面（如 SDF / CSG）在几何求交、渲染投射与拓扑编辑上的优劣与计算特性。
    \item \textbf{[开发]} 通过多边形表面曲面细分与半边拓扑缝合，独立构建基础多面体几何模型。
    \item \textbf{[开发]} 编程实现移动立方体算法，从三维隐式等值面场或体素场中提取闭合的连续三角多边形网格。
    \item \textbf{[开发]} 运用过程式规则或分形算法（如分形布朗运动 fBm 或细分菱形-方形算法），程序化生成高起伏逼真的分形山脉地形。
    \item \textbf{[开发]} 从工业级激光三维扫描仪采集的散乱无序空间点云中，编程实现泊松表面重构算法生成光滑闭合网格。
    \item \textbf{[应用]} 基于基础几何图元（如单位立方体、二次曲面球体与圆柱），通过 CSG 树状布尔交并差运算构建复杂机械零件模型。
    \item \textbf{[评估]} 全面对比多边形网格、隐式距离场与体素体数据在内存消耗、几何拓扑修改延迟与最终渲染光照质感上的技术权衡。
\end{enumerate}

\subsection{高级着色与进阶渲染}
\textbf{【学时建议】}：6 领域核心学时。

本单元深入探究真实物理光线传输机制的精确求解算法、微表面微观物理材质建模以及现代图形处理器（GPU）硬件架构的性能极值优化。

\subsubsection{领域核心主题}
\begin{itemize}
    \item \textbf{渲染方程的高阶数值求解与近似算法}：
    \begin{itemize}
        \item 分布式光线追踪与纯蒙特卡洛积分收敛原理。
        \item 经典单向蒙特卡洛路径追踪算法与重要性采样。
        \item 光子映射技术（Photon Mapping，处理焦散效应 Caustics 与漫反射间接光照）。
        \item 双向路径追踪 (BDPT)。
        \item 都会光线传输算法（Metropolis Light Transport, MLT，基于马尔可夫链蒙特卡洛 MCMC 状态空间探索）。
    \end{itemize}
    \item \textbf{真实物理成像维度与时间/镜头集成}：
    \begin{itemize}
        \item 时间域采样：相机快门开启过程中的时间积分与真实动态模糊。
        \item 空间光学透镜系统：有限薄透镜孔径成像、景深与对焦控制。
        \item 连续光学光谱频率：全光谱光传输与色散棱镜彩虹物理效果。
    \end{itemize}
    \item \textbf{高级阴影计算体系}：阴影贴图技术、采样自走样消除、透视阴影修正、级联阴影贴图 (CSM)与方差阴影贴图 (VSM)。
    \item \textbf{几何遮挡剔除技术}：硬件遮挡查询、潜在可见集（PVS）与软件光栅化剔除。
    \item \textbf{基于物理的双向散射分布函数 (BSDF) 与微表面理论}：
    \begin{itemize}
        \item 双向反射分布函数（BRDF）与双向透射分布函数（BTDF）。
        \item 经典微表面理论：法线分布函数（NDF，如 GGX / Trowbridge-Reitz）、几何遮蔽函数（Smith 几何项）与菲涅尔项（Schlick 近似公式）。
    \end{itemize}
    \item \textbf{次表面散射物理模拟}：偶极子扩散近似模型及其在皮肤、蜡质、玉石等半透明多孔介质渲染中的应用。
    \item \textbf{真实多面体与多边形区域面光源采样}：多边形光源的解析解（如基于 LTC 线性变换余弦）与软阴影半影区物理过渡。
    \item \textbf{层次化深度缓冲加速}：GPU 硬件 Early-Z 早期深度剔除流水线。
    \item \textbf{光场理论与基于图像的渲染技术}：全光函数、视差贴图、环境光照映射（IBL）与神经辐射场呈现机理。
    \item \textbf{非真实感艺术风格渲染}：卡通着色、轮廓边缘检测（Sobel 与法线/深度不连续性提取）、手绘风格与素描笔触合成。
    \item \textbf{现代 GPU 高性能计算与微架构解剖}：流式多处理器（SM）组织、Warp 线程束调度、显存访存合并、硬件光线追踪加速内核（如 RT Core BVH 遍历与求交引擎）。
    \item \textbf{人类视觉感知在高阶着色中的映射}：明暗适应性动态调节、人眼对光强高频噪声的掩蔽效应、频闪融合与高帧率刷新对真实感的塑造。
\end{itemize}

\subsubsection{示范性学习成果与能力目标}
\begin{enumerate}
    \item \textbf{[证明]} 通过数学推导，证明蒙特卡洛路径追踪算法如何通过对全场景表面立体角的随机光线采样，逼近渲染方程的真实解析解。
    \item \textbf{[证明]} 分析并严格证明某种渲染数值积分算法的理论属性（例如：完全性、统计一致性 Consistent 与无偏性 Unbiased）。
    \item \textbf{[分析]} 定量分析一个典型片元着色算法（如带软阴影过滤的 PBR 着色）在特定分辨率下的显存带宽消耗与浮点计算吞吐量需求。
    \item \textbf{[开发]} 基于现代实时渲染 API，独立编程实现一种非平凡的高阶着色算法（如卡通着色结合轮廓提取、四阶级联阴影贴图 CSM 或实时环境光遮蔽 SSAO）。
    \item \textbf{[应用]} 说明特定传统绘画艺术技法（如明暗对照法 Chiaroscuro、水彩晕染效果）如何在计算机着色器算法中得到精准的物理与数学映射。
    \item \textbf{[分析]} 给定一幅高拟真计算机合成图像，准确反向分析识别其在生成时所运用的特定核心图形学着色技术（如次表面散射、菲涅尔反射、全局间接焦散）。
    \item \textbf{[开发]} 脱离任何现代上层图形 API，在最基础的逐像素缓冲区纯 CPU 内存层面，手工编写代码实现光栅化与光照着色。
    \item \textbf{[开发]} 编写一个支持任意三角形三维场景的递归光线追踪渲染器，完整支持经典 Phong/Blinn-Phong BRDF 模型、理想镜面高阶反射与折射光线跟踪。
\end{enumerate}

\subsection{计算机动画}
\textbf{【学时建议】}：6 领域核心学时。

本单元涵盖二维与三维虚拟对象随时间演化发生几何形变、骨骼驱动与空间运动的核心数学规律、算法实现及艺术表现原则。

\subsubsection{领域核心主题}
\begin{itemize}
    \item \textbf{经典动画十二基本法则}：
    \begin{itemize}
        \item 挤压与拉伸（Squash and Stretch，维持体积恒定下的弹性形变）。
        \item 时间节奏与运动节拍。
        \item 预备动作（Anticipation，唤起受众对后续主要动作的心理预期）。
        \item 演出布局与视线引导。
        \item 跟随动作与重叠动作（Follow Through and Overlapping Action，动量滞后传递）。
        \item 逐帧绘制与原画关键帧推进。
        \item 慢入与慢出（Slow In and Slow Out，加减速物理缓冲惯性）。
        \item 弧线运动轨迹（Arcs，自然生物关节转动的几何规律）。
        \item 夸张手法与视觉感染力。
    \end{itemize}
    \item \textbf{关键帧动画技术与插值数学}：
    \begin{itemize}
        \item 关键帧标定与时序曲线编辑。
        \item 经典插值数学方法：线性插值、四元数球面线性插值、三次样条插值与贝塞尔加速度平滑曲线。
    \end{itemize}
    \item \textbf{正向运动学与逆向运动学}：
    \begin{itemize}
        \item 正向运动学 (FK)：骨骼层级树状结构、局部变换矩阵沿父子节点的向下递归累乘。
        \item 逆向运动学 (IK)：末端执行器约束求解、循环坐标下降法（CCD）、FABRIK 算法以及基于雅可比矩阵伪逆的数值优化求解。
    \end{itemize}
    \item \textbf{骨骼蒙皮算法族}：
    \begin{itemize}
        \item 蒙皮权重分配与顶点绑定。
        \item 线性混合蒙皮 (LBS)及其固有的关节点“糖果纸包装扭曲”与体积丢失缺陷。
        \item 双四元数蒙皮 (DQS)在刚性体积保持上的数学优势。
    \end{itemize}
    \item \textbf{骨骼角色装配绑定体系}：骨骼层级约束、反向控制器、极向量约束。
    \item \textbf{混合变形与表情系统}：Blendshapes（形态键/融合成形）、姿态空间形变（Pose Space Deformation, PSD，用于矫正肌肉挤压拉伸隆起）。
    \item \textbf{空间几何变换系统与旋转表征体系}：
    \begin{itemize}
        \item 三维平移与各向异性缩放/剪切。
        \item 三维空间旋转数学表征之比较：欧拉角及其固有的万向节死锁现象；四元数（Quaternions，单位四元数乘法与旋转向量旋转）；轴角表示法与李代数指数映射。
        \item 变换矩阵的链式乘法顺序对最终姿态的决定性影响（SRT 顺序 vs XYZ 轴顺序）。
    \end{itemize}
    \item \textbf{虚拟摄影机轨迹动画}：观察注视点（Look-at 矩阵构建与时序插值跟踪）、变焦与景深对焦动画。
    \item \textbf{动作捕捉技术}：
    \begin{itemize}
        \item 光学反光标记点追踪、惯性传感器（IMU）追踪系统设置与数据标定管线。
        \item 动作捕捉伦理与社会规范：捕捉数据的隐私知情权；坚决摒弃“默认通用行走步态”的狭隘工程设定——现实世界中不存在所谓“标准人类步态”，系统设计必须兼顾残障、老龄与多元生理群体的非标运动形态。
    \end{itemize}
\end{itemize}

\subsubsection{示范性学习成果与能力目标}
\begin{enumerate}
    \item \textbf{[计算]} 给定骨骼层级树与各关节点的局部旋转矩阵，利用正向运动学（FK）自顶向下递归计算指定肢体关节点的世界绝对空间坐标与朝向。
    \item \textbf{[计算]} 针对给定的肢体末端执行器目标空间位置与朝向，采用逆向运动学（IK）算法求解出各关节所需配置的角度姿态。
    \item \textbf{[对比]} 深入对比欧拉角、旋转矩阵、四元数与轴角在存储紧凑度、插值连续平滑性以及避免奇异性（万向节死锁）上的核心权衡。
    \item \textbf{[开发]} 编程实现三次样条插值或四元数球面线性插值算法，平滑生成连续的关键帧中间姿态与运动过渡。
    \item \textbf{[应用]} 熟练运用工业级动画建模软件（如 Blender 或 Maya），结合骨骼装配体系与隐式电位球或多边形网格，完成一个有机生命角色的基础动作装配与行走动画序列制作。
\end{enumerate}

\subsection{计算机物理仿真}
\textbf{【学时建议】}：6 领域核心学时。

计算机物理仿真在数理根基上与计算科学紧密相连。但在计算机图形学与交互技术语境下，仿真技术被赋予了截然不同的学科使命：它不再片面追求严苛的高精度物理预测模型，而是追求“物理真实感与艺术家创作意图之间的精妙平衡”。例如，科学计算关注“精确建模流体的微观表面张力偏微分方程”，而图形学则要求“受控地生成极具视觉美感与冲击力的皇冠状跌落水花”。

\subsubsection{领域核心主题}
\begin{itemize}
    \item \textbf{碰撞检测与动态碰撞响应 (Collision Detection \& Response)}：
    \begin{itemize}
        \item 基于有符号距离场（SDF）的快速内外穿透判断与梯度法向量反弹计算。
        \item 球-球相交、点-三角形空间包含判定、边-边最小距离分离轴测试。
        \item 连续碰撞检测（CCD）规避高速物体运动“穿模隧道效应”。
    \end{itemize}
    \item \textbf{基于随机噪声的过程式运动动画}：Perlin 噪声、Simplex 噪声在风场扰动与复杂海浪表面起伏中的程序化叠加。
    \item \textbf{粒子动力学系统}：
    \begin{itemize}
        \item 时间连续常微分方程的离散数值积分器：前向欧拉法（Forward Euler，条件不稳定）、中点法、蛙跳积分法、辛欧拉法与 Verlet / Velocity-Verlet 积分。
        \item 质点-弹簧阻尼网络（Mass-Spring Networks，用于布料织物与柔软弹性形变模拟）。
        \item 基于位置的动力学（Position-Based Dynamics, PBD，无条件数值稳定、快速约束投影求解）。
        \item 规则驱动的群体仿生学系统：Boids 经典鸟群三法则（分离、对齐、凝聚）与高密度人群疏散仿真。
        \item 刚体运动学与动力学（冲量积分法与角动量牛顿-欧拉方程）。
    \end{itemize}
    \item \textbf{基于欧拉网格的流体动力学}：
    \begin{itemize}
        \item 不可压缩流体纳维-斯托克斯方程。
        \item 半拉格朗日平流算法（Semi-Lagrangian Advection，无条件稳定的一阶回溯插值）。
        \item 亥姆霍兹-霍奇分解与泊松压力投影方程求解器（Pressure Projection，保持速度场无散度 $\nabla \cdot \mathbf{u} = 0$）。
    \end{itemize}
    \item \textbf{高度场与地形演化系统}：
    \begin{itemize}
        \item 地表沉积物输运、水力侵蚀与热风化侵蚀模型。
        \item 水波仿真：波动方程与浅水方程 (SWE)。
    \end{itemize}
    \item \textbf{基于离散规则的复杂系统}：L-系统（L-Systems，植物形态生长）、空间占优殖民系统（Space-colonizing Systems，叶脉与树枝向光性扩展）、康威生命游戏（Game of Life，元胞自动机）。
    \item \textbf{仿真必备高级几何数据结构}：
    \begin{itemize}
        \item 密集三维体素标量场/矢量场。
        \item 层次化稀疏体素场（Sparse Volumes，如 OpenVDB，有效节约烟雾火焰的巨大空白内存开销）。
        \item 自适应网格细分场（Adaptive Volumes，在激波与飞溅界面动态加密）。
        \item 携带多物理属性的粒子点集（Points with Attributes，质量、速度、密度、相态）。
        \item 非结构化三角形汤。
        \item 二维高程图。
    \end{itemize}
    \item \textbf{目标导向的仿真技术选型映射体系}：
    \begin{itemize}
        \item 粒子系统应用：粒子数值积分（前向欧拉、中点法、蛙跳法）。
        \item 刚体动力学：质点系统扩展、三角形-顶点/边-边相交碰撞检测。
        \item 织物布料模拟：基于质点-弹簧模型或 PBD 位置约束、自碰撞检测机制。
        \item 基于粒子的水体流体：光滑粒子流体动力学 (SPH)核函数插值、基于 SDF 的刚度障碍物碰撞反弹。
        \item 基于网格的烟雾与火焰：半拉格朗日平流、交错网格 MAC 离散化、泊松压力投影。
        \item 混合粒子-网格水体模拟：FLIP（Fluid-Implicit-Particle）与 PIC（Particle-In-Cell）混合算法。
    \end{itemize}
\end{itemize}

\subsubsection{示范性学习成果与能力目标}
\begin{enumerate}
    \item \textbf{[开发]} 基于经典牛顿力学定律，编程实现粒子动力学物理仿真引擎，支持辛欧拉法、中点法或 Verlet 积分器，完成诸如弹性蠕虫链、弹性摆球等物理现象的交互仿真。
    \item \textbf{[对比]} 深入对比不同流体仿真技术（如纯拉格朗日 SPH 粒子法、欧拉网格半拉格朗日法与混合 FLIP 算法）在表现弹道式飞溅水花、爆炸烟尘与高扩散火光时的物理直观性、数值耗散与算力开销。
    \item \textbf{[开发]} 独立设计并编程实现一个支持实时交互的二维/三维烟雾求解器，支持用户实时添加障碍物、扰动速度场并注入热源与烟尘浓度。
\end{enumerate}

\subsection{沉浸式技术}
\textbf{【学时建议】}：6 领域核心学时。

沉浸式技术重新定义了人与虚拟数字空间的交互边界。本单元涵盖虚拟现实（VR）、增强现实（AR）与混合现实（MR）所依托的显示光学、空间立体追踪、超低时延渲染管线、人因工程及空间隐私伦理。

\subsubsection{领域核心主题}
\begin{itemize}
    \item \textbf{扩展现实概念谱系辨析}：虚拟现实（VR）、增强现实（AR）、混合现实（MR）与扩展现实（XR）的技术边界；Milgram 提出的“现实-虚拟连续统”理论。
    \item \textbf{双目立体视觉显示技术}：左右眼视差生成、微显示器光机、菲涅尔透镜与薄煎饼透镜光学畸变预校正、色差补偿、眼动辐辏调节冲突 (VAC)。
    \item \textbf{观察者空间位姿追踪技术}：
    \begin{itemize}
        \item 追踪架构：由外向内追踪（Outside-In，利用外部红外基站或光捕相机）与由内向外追踪（Inside-Out，利用头显机载摄像机与 VIO 视觉惯导 SLAM 算法）。
        \item 空间 6 自由度（6-DoF）追踪维度：头部位置与朝向、双手控制器与裸手手势追踪、全身骨骼姿态捕捉与视线眼动追踪。
    \end{itemize}
    \item \textbf{沉浸式超低延迟渲染与可见性优化}：
    \begin{itemize}
        \item 动光延迟（Motion-to-Photon Latency，从用户头部物理转动到显示屏光子完全刷新出对应画面的物理延迟，须严格压制在 20ms 以内以防止生理眩晕）。
        \item 异步时间扭曲（ATW）、异步空间扭曲（ASW）与重投影技术。
        \item 注视点渲染（Foveated Rendering，利用眼动追踪在人眼中心凹视野区域提供全高清画质，周边外围视野大幅降采样）。
        \item 基于图像的全景 VR 系统与多细节层次（LOD）动态调度。
    \end{itemize}
    \item \textbf{网络化分布式虚拟现实}：多用户跨局域网/广域网的虚拟协同空间、空间状态同步、网络抖动补偿与空间音频定位渲染。
    \item \textbf{虚拟现实中的交互式三维空间建模}。
    \item \textbf{沉浸式技术行业深度赋能}：精准数字化外科手术导航、复杂战术对抗仿真推演、高端装备制造维修全息装配指引、沉浸式科学数据多维漫游。
    \item \textbf{沉浸式系统的人因安全与物理风险}：
    \begin{itemize}
        \item 晕动症/眩晕病的神经病理机制（前庭感官感知运动与视网膜视觉静止之间的感觉冲突）。
        \item 物理环境遮蔽安全隐患：沉浸式 VR 完全阻断现实视觉，导致室内跌倒摔伤、撞击实体家具与墙壁风险。
    \end{itemize}
    \item \textbf{沉浸式技术的无障碍普惠设计}：针对严重肢体运动障碍用户的注视交互适配、针对卧床不起病患的空间视锥倾斜补偿。
    \item \textbf{空间计算隐私与伦理风险 (交叉参考 SEP)}：
    \begin{itemize}
        \item 个人物理私密空间数字重构的泄露风险（头显机载相机对私人卧室、办公室空间的连续深度点云扫描与环境建模数据上传云端）。
        \item 生物特征数据滥用（步态、瞳孔缩放、身体骨骼尺寸比例）。
        \item 视觉强沉浸感导致的心理成瘾、现实脱节以及沉浸过程中用户对周边物理攻击的防御丧失。
    \end{itemize}
\end{itemize}

\subsubsection{示范性学习成果与能力目标}
\begin{enumerate}
    \item \textbf{[开发]} 编程渲染生成一对符合人类平均瞳距（IPD）几何视差的左/右眼立体匹配图像对。
    \item \textbf{[对比]} 详细阐明并对比由外向内与由内向外空间定位追踪技术在部署成本、追踪范围、遮挡敏感性与计算负荷上的核心优缺点。
    \item \textbf{[对比]} 深入对比基于纯几何建模的虚拟现实系统与基于全景立体图像重构的虚拟现实系统在自由漫游度与计算资源需求上的本质差异。
    \item \textbf{[评估]} 在高并发网络化多人协作沉浸式系统中，评估并设计保证动作同步实时性与多端数据状态强一致性的网络传输方案。
    \item \textbf{[设计]} 面向医院洁净手术室的心脏介入手术场景，撰写一套详尽的增强现实（AR）辅助导航应用程序功能需求与软硬件技术规范。
    \item \textbf{[评估]} 依据人机交互（HCI）与无障碍通用设计准则，对一款现存的沉浸式 XR 应用进行全维度的无障碍包容性评估并出具改进报告。
    \item \textbf{[识别]} 明确指出引发虚拟现实晕动病的核心关键软硬件技术指标（如动光延迟 MTP、双目视差校准度、画面帧率抖动），并论述抑制眩晕的系统工程优化手段。
\end{enumerate}

\subsection{交互技术}
\textbf{【学时建议】}：6 领域核心学时。

现代交互式计算机图形学构成了从桌面排版工具到大型多人在线游戏、元宇宙沉浸空间等所有高实时性交互软件的运转灵魂。本单元重点培养学生跨越底层硬件中断、事件驱动消息机制直至多模态感官交互全链路的系统构建能力。

\subsubsection{领域核心主题}
\begin{itemize}
    \item \textbf{事件驱动编程范式}：
    \begin{itemize}
        \item 鼠标指针移动、滚轮与物理触控点击事件。
        \item 键盘矩阵按键按下与释放状态机。
        \item 语音识别输入流解析与自然语言交互触发。
        \item 物理环境传感器数据触发（光感、陀螺仪、接近传感器）。
        \item 跨组件与跨线程的消息传递通信。
        \item 异步网络套接字事件轮询（如 Select/Epoll 事件模型）。
        \item 底层硬件中断事件处理与驱动级回调机制。
    \end{itemize}
    \item \textbf{单通道经典图形用户界面 (WIMP 范式)}：视窗、图标、菜单、定点设备的状态变迁与事件分发路由机制。
    \item \textbf{手势交互接口}：多点触控手势识别、空中自由手势几何匹配（如捏合、滑动）、无障碍替代方案设计（如为手部缺失或运动神经损伤者提供软件逻辑层面的“虚拟点赞”事件注入）。
    \item \textbf{触觉与力反馈接口}：
    \begin{itemize}
        \item 外部线性振动执行器（LRA 与 ERM 电机）。
        \item 具备压力与纹理触觉感知的力反馈数据手套。
        \item 限制关节运动的大负载力反馈机械外骨骼。
    \end{itemize}
    \item \textbf{多模态融合交互系统}：视线追踪、语音指令、手势捕捉与微表情识别的时空融合与多义性消歧。
    \item \textbf{基于人工智能的前沿沉浸式交互接口}：
    \begin{itemize}
        \item 脑机接口（BCI）：基于非侵入式脑电图（EEG）电极阵列的稳态视觉诱发电位（SSVEP）与运动想象意图解码。
        \item 集成眼动追踪的混合现实智能头显：凝视选择与眼部眨眼触发。
        \item 智能全息光波导 AR 眼镜的极简微交互。
    \end{itemize}
    \item \textbf{无障碍设计准则与通用包容性设计 (交叉参考 SEP)}：为肢体、听觉、认知等各类障碍群体提供无障碍等价交互通道。
\end{itemize}

\subsubsection{示范性学习成果与能力目标}
\begin{enumerate}
    \item \textbf{[开发]} 编写一个完整的二维图形小游戏，能够对单通道鼠标移动、点击以及键盘状态事件做出即时、平滑的逻辑响应。
    \item \textbf{[开发]} 编写嵌入式控制程序，使单片机电路能够根据外接可变电位器（旋转电阻）的模拟阻值变动，平滑调节虚拟图形对象的多维属性。
    \item \textbf{[开发]} 面向移动智能终端，开发一个能精准响应单指滑动、双指缩放及多点复杂手势的交互式移动应用程序。
    \item \textbf{[应用]} 利用计算机视觉库或手势追踪外设，编写程序实现通过手部空中手势对三维场景中虚拟物体的空间抓取与平移旋转。
    \item \textbf{[设计]} 设计并编程实现一个集成立体触觉反馈（如变频振动或力阻尼）的交互应用程序，向用户传递虚拟物体的表面粗糙度与刚度质感。
    \item \textbf{[开发]} 架构并实现一个复杂的响应式图形交互系统，使其能够稳定并发处理来自网络套接字、定时器硬件中断以及异步用户外设的多源事件触发器。
\end{enumerate}

\subsection{数字图像处理}
\textbf{【学时建议】}：6 领域核心学时。

数字图像处理涵盖面向多种工程目标的图像分析、滤波、增强、编码与几何变换，其核心主旨在于提升视觉感知品质与实现图像内容的智能操纵。该单元是计算机视觉（CV）不可动摇的技术基石。

\subsubsection{领域核心主题}
\begin{itemize}
    \item \textbf{二值数学形态学处理}：
    \begin{itemize}
        \item 像素连通分量标记（4-连通与 8-连通域分析）。
        \item 经典形态学算子：膨胀、腐蚀、开运算与闭运算。
        \item 连通区域几何属性统计提取：面积、周长、质心坐标、方向角、主轴长度与圆度。
    \end{itemize}
    \item \textbf{色彩直方图分析与对比度操纵}：直方图统计表征模型、直方图均衡化、局部自适应直方图均衡化（CLAHE）与动态范围归一化拉伸。
    \item \textbf{图像空域滤波与图像增强}：
    \begin{itemize}
        \item 二维离散空间卷积运算原理。
        \item 图像平滑与模糊算子：均值滤波、高斯低通滤波、双边滤波（Bilateral Filter，保边去噪）。
        \item 图像锐化算子：拉普拉斯算子、Sobel 梯度算子、反锐化掩模技术。
        \item 频域滤波技术：二维快速傅里叶变换（FFT）、理想低通/高通滤波器、巴特沃斯滤波器与高斯频域滤波。
    \end{itemize}
    \item \textbf{图像退化物理建模与图像复原}：加性高斯白噪声、椒盐噪声、运动模糊退化函数模型；逆滤波、维纳滤波；现代图像修复技术（Inpainting，基于偏微分方程与泊松图像编辑 Poisson Image Editing）。
    \item \textbf{数字图像编码与信息压缩理论}：
    \begin{itemize}
        \item 空间冗余、时间冗余与人眼视觉心理冗余。
        \item 无损熵编码：香农-哈夫曼编码（Huffman Coding）、算术编码与行程编码（RLE）。
        \item 正交变换编码：离散余弦变换（DCT，JPEG 核心）、离散小波变换（DWT，JPEG 2000 核心）。
        \item 奈奎斯特-香农采样定理与频谱混叠规避。
        \item 数字盲水印嵌入与数字版权管理（DRM）、抗几何攻击鲁棒性与隐蔽性伦理考量。
    \end{itemize}
    \item \textbf{前沿深度学习视觉技术}：卷积神经网络（CNN，特征提取层、池化层与残差连接）、视觉 Transformer 架构（ViT，自注意力机制在图像分块表征中的应用）、生成对抗网络（GAN）与前沿扩散去噪概率模型。
    \item \textbf{伦理、司法防伪与社会公平挑战 (SEP)}：换脸与深度伪造（Deepfakes）的自动取证与鉴伪技术；计算机视觉算法因训练样本失衡而在特定肤色、发型或少数族裔群体上出现的严重误识别、司法误判与社会歧视隐患。
\end{itemize}

\subsubsection{示范性学习成果与能力目标}
\begin{enumerate}
    \item \textbf{[应用]} 运用形态学膨胀与腐蚀算子，对带有毛刺噪声的二值图像进行平滑与孔洞填补，并准确分离粘连的微小图元。
    \item \textbf{[开发]} 编程提取图像的 HSV 色彩空间矩阵，在独立保持亮度与饱和度不变的前提下，操纵调整图像的色相通道。
    \item \textbf{[开发]} 编写程序实现高通滤波算子（在空间域利用卷积核，或在频域利用傅里叶高通掩模），实现对图像微弱边缘轮廓的高对比度增强。
    \item \textbf{[开发]} 采用高级图像修复补全算法（如泊松图像无缝融合 Poisson Editing 或基于样本块的修补技术），对被严重损毁或人为涂抹的图片进行无痕修复还原。
    \item \textbf{[开发]} 编程对二维图像进行傅里叶正变换，通过在频域平面内有选择性地滤除特定周期性高频条纹噪声，实现图像质量逆向增强。
    \item \textbf{[评估]} 深入阐明人脸识别算法在特定应用场景下的技术缺陷与伦理陷阱；针对“人脸识别能否以及如何做到合乎工程伦理的应用”撰写深度学术评述报告。
\end{enumerate}

\subsection{实体与物理计算}
\textbf{【学时建议】}：6 领域核心学时。

本单元与专用平台开发（SPD）以及人机交互（HCI）高度交织，旨在打破传统计算机“屏幕-键鼠”的物理禁锢，将计算能力拓展至真实物理世界的实体对象、智能空间、传感器网络与数字制造链条中。

\subsubsection{领域核心主题}
\begin{itemize}
    \item \textbf{与真实物理世界的通信与交互感知}：
    \begin{itemize}
        \item 物理环境数据感知采集：模拟信号与模拟-数字转换（ADC）、温度/光敏电阻/超声波测距/红外接近/加速度与陀螺仪传感器数据流处理。
        \item 物理世界驱动控制：脉冲宽度调制（PWM）技术驱动舵机、直流电机驱动器、高精度双极性步进电机、电磁阀与继电器。
    \end{itemize}
    \item \textbf{微控制器的事件驱动固件架构}：定时器硬件中断、外部引脚电平跳变中断服务函数（ISR）、环形串口缓冲队列与非阻塞式主循环调度。
    \item \textbf{与物理实体的制造连接}：
    \begin{itemize}
        \item 计算机辅助设计（CAD）：实体造型、参数化约束、装配体干涉检查。
        \item 计算机辅助制造（CAM）：刀具路径规划、G 代码与 M 代码生成。
        \item 现代数字制造工艺路线：
        \begin{itemize}
            \item 人机交互快速原型设计（HCI Prototyping，低成本纸板与激光亚克力快速成型）。
            \item 增材制造技术（3D Printing）：熔融沉积成型（FDM）、立体光固化成型（SLA）、选择性激光烧结（SLS）工艺特性对比。
            \item 减材制造技术：数控机床铣削、高精度激光切割与水刀切割。
            \item 成型制造工艺：真空热吸塑成型、金属板料冲压折弯。
        \end{itemize}
    \end{itemize}
    \item \textbf{物联网感知与互联体系}：
    \begin{itemize}
        \item 网络互联拓扑结构与轻量级协议栈（MQTT、CoAP）。
        \item 短程与低功耗无线通信标准：Wi-Fi (802.11)、低功耗蓝牙 (BLE)、Zigbee、LoRa 远程扩频通信与近场通信 (RFID/NFC)。
    \end{itemize}
    \item \textbf{实体计算安全与隐私伦理 (SEP)}：物联网终端物理逆向工程漏洞、家用常开麦克风与环境传感器的隐私嗅探防范、固件空中升级（OTA）防篡改机制。
\end{itemize}

\subsubsection{示范性学习成果与能力目标}
\begin{enumerate}
    \item \textbf{[开发]} 构建一个基础硬件控制电路（包含微控制器、限流电阻与物理按键开关），编写固件实现通过按键消抖触发点亮或关闭高亮发光二极管（LED）。
    \item \textbf{[开发]} 搭建并联调一个闭环实体机电控制系统，使舵机能够根据超声波或红外距离传感器的实时测距数值，做出精准角度旋转跟踪。
    \item \textbf{[开发]} 利用光敏传感器采集环境照度，通过编程动态驱动外部物理执行器（如平滑调节 LED 阵列的发光亮度或驱动显示屏幕中图形的虚拟色彩）。
    \item \textbf{[开发]} 在专业计算机辅助设计（CAD）工业软件中独立构建一个复杂的三维实体零件，达成以下指标：
    \begin{itemize}
        \item 在 CAD 系统中演示并论证空间仿射变换（平移、旋转、镜像）的底层数学实现。
        \item 演示对象的实例化复用机制。
        \item 制定完整的零件数字化制造工序计划，详细评估材料损耗与加工工时成本。
        \item 实际操控 3D 打印机或 CNC 机床完成实体加工，对比实测物理零件与理论设计图纸的公差差异，剖析加工误差来源。
    \end{itemize}
    \item \textbf{[开发]} 手工编写标准数控 G 代码与 M 代码，用于控制数控铣床自动雕刻一个具有特定拓扑连通性的三维迷宫零件，并利用 CAM 仿真软件检验刀路碰撞。
    \item \textbf{[评估]} 面向适老化智能医疗药盒的研发任务，系统论证应选用以太网、Wi-Fi、BLE、RFID/NFC 或蜂窝网络中的哪种方案实现云端互联，并独立制作可运行的硬件原型系统。
    \item \textbf{[对比]} 深入对比增材制造、减材制造与吸塑成型工艺在材料强度、几何自由度、单件打样成本与规模化量产时的经济技术适用边界。
\end{enumerate}

\section{专业素养与思维倾向}

在图形学与交互技术的系统研习中，必须高度注重塑造学生以下几维关键的工程师素养与思维倾向：
\begin{itemize}
    \item \textbf{积极主动与自发探索}：面对未知的底层图形 API、晦涩的物理着色方程与复杂的物理外设驱动，具有自主查阅前沿文献与工业规范的主动探索精神。
    \item \textbf{自我驱动与终身学习}：图形硬件与实时渲染技术日新月异，从业者需保持高度的自驱力，能够紧跟全球学术界（如 SIGGRAPH）与工业界引擎前沿演进。
    \item \textbf{目标导向与结果思维}：在算力、显存带宽与物理帧率强约束下，始终以“满足真实感与帧率要求”为最终导向开展架构设计。
    \item \textbf{职业操守与伦理自律}：在处理涉及人脸生物特征、空间点云扫描、数字版权与生成式深度伪造内容时，恪守工程道德，尊重知识产权，严格保护用户个人隐私。
    \item \textbf{敏捷适应与灵活变通}：面对跨平台图形后端（DirectX、Vulkan、Metal、WebGPU）以及异构计算环境，具备极强的代码架构迁移与兼容变通能力。
    \item \textbf{团队协作与跨界融合}：图形工程通常处于软件工程师、数学算法专家、三维数字艺术家、工业设计师以及产品经理的交叉枢纽，必须具备强大的跨学科团队协作素养。
    \item \textbf{敏锐响应用户反馈}：高度关注真实用户的感官体验、交互延迟以及生理不适，能敏锐依据用户体验数据与可用性测试迭代优化软件。
    \item \textbf{精益求精与严谨细致}：深入推敲每一个像素着色计算精度、浮点数舍入误差、网格法线奇异点以及内存缓冲区边界溢出，防范图形破面与渲染穿模。
    \item \textbf{突破定势与开拓创新}：面对极其耗时的物理级光线模拟计算，敢于跳出传统蛮力计算框架，创造性地运用微表面理论、降噪神经网络或启发式空间剖分技术进行性能突围。
    \item \textbf{勇于担当与后果意识}：深刻意识到图形交互算法对物理现实安全的直接影响（如自动驾驶机器视觉盲区、医疗 AR 手术导航定位偏差），对软件系统的一切潜在后果负有敬畏与担当之心。
    \item \textbf{规范的代码审查能力}：善于开展深入的同行代码审查，既能发现他人代码中的显存泄露与性能瓶颈，亦能以开放的心态吸纳同行改进建议。
    \item \textbf{多维立体的高效沟通素养}：
    \begin{itemize}
        \item \textbf{口头表达}：能面向非技术背景的艺术设计师清晰阐释复杂着色器技术参数的物理内涵。
        \item \textbf{书面表达}：能撰写结构严密、数学定义准确的技术白皮书、系统架构设计说明书与学术研究报告。
        \item \textbf{代码表达}：能编写结构优雅、自解释性强、注释规范且符合工业级编码风格规范的高质量代码。
    \end{itemize}
\end{itemize}

\section{数学与基础数据结构要求}

图形学与交互技术建立在扎实完备的现代应用数学基石之上：

\subsection{必须具备的数学核心}
\begin{itemize}
    \item \textbf{线性代数}：
    \begin{itemize}
        \item 空间点、笛卡尔坐标系、仿射空间与四维齐次坐标系统。
        \item 向量代数：向量加减、标量乘法、点积（Dot Product，用于夹角与漫反射余弦计算）、叉积（Cross Product，用于表面法向量求解与三角形朝向判定）。
        \item 矩阵代数：矩阵乘法、逆矩阵、转置矩阵以及伴随矩阵（用于法向量变换逆转置矩阵推导）。
        \item 仿射变换矩阵构造：平移、旋转、非均匀缩放、剪切与透视投影矩阵全流程推导。
    \end{itemize}
    \item \textbf{微积分}：标量函数与多元空间矢量函数的连续性概念、偏导数、空间梯度、定向导数与多重微积分基本概念。
\end{itemize}

\subsection{建议掌握的进阶数学}
\begin{itemize}
    \item \textbf{高阶线性代数}：
    \begin{itemize}
        \item 矩阵特征值与特征向量分解（Eigen Decomposition，用于主成分分析 PCA 与网格主曲率计算）。
        \item 高斯消元法与矩阵三角分解（LU 分解，用于刚体与流体线性泊松方程组快速求解）。
        \item 奇异值分解（SVD，用于三维刚性点云精确配准与最佳位姿拟合）。
    \end{itemize}
    \item \textbf{进阶代数与数值微积分}：
    \begin{itemize}
        \item 四元数代数体系（Quaternions，哈密顿算子、四元数插值 Slerp 与万向节死锁彻底消除）。
        \item 常微分方程数值解法与蒙特卡洛积分方法。
    \end{itemize}
\end{itemize}

\subsection{核心数据结构要求}
\begin{itemize}
    \item \textbf{必须熟练掌握的底层数据结构}：
    \begin{itemize}
        \item 有向无环图（DAG）：用于层次化场景图、着色器材质节点图。
        \item 定长元组：空间几何点、三维物理向量、固定阶数齐次变换矩阵。
        \item 密集多维张量数组（Dense 1D, 2D, 3D Arrays）：帧缓冲区、深度缓冲区、二维纹理采样图元、三维体素标量场。
    \end{itemize}
    \item \textbf{极具工程价值的进阶数据结构}：
    \begin{itemize}
        \item 结构体数组与数组结构体的底层访存布局对偶性：理解数据导向型设计（DOD）在 CPU 缓存局部性与 GPU SIMD 向量化并行中的极致性能影响。
        \item 高阶空间检索树结构：$K$-d 树、四叉树、八叉树、包围盒层次结构（BVH）以及用于图像压缩的哈夫曼树。
    \end{itemize}
\end{itemize}

\section{交叉主题与跨学科共享概念}

计算机图形学与交互技术并非孤立存在，它与计算机科学各分支学科紧密交融：
\begin{itemize}
    \item \textbf{沉浸式技术 (GIT) 与人机交互 (HCI)}：空间三维用户体验、凝视交互、人机认知负荷评估。
    \item \textbf{交互技术 (GIT) 与 人机交互/GUI 编程以及计算科学 (CN) 交互可视化}：复杂事件驱动模型、状态图迁移管理。
    \item \textbf{过程式生长建模 (GIT) 与编程语言理论 (FPL)}：L-系统文法与巴科斯-诺尔范式（BNF）形式文法的对偶性——利用文法规则派生三维几何分形图像。
    \item \textbf{物理仿真 (GIT) 与计算科学 (CN)}：连续动力学方程离散化与数值阻尼控制。
    \item \textbf{数据可视化 (GIT) 与人工智能 (AI) 及专用平台开发 (SPD)}：高维机器学习特征流形嵌入可视化分析。
    \item \textbf{数字图像处理 (GIT) 与 人工智能/机器视觉}：卷积特征提取、图像分割、双目立体匹配。
    \item \textbf{实体计算与计算机动画 (GIT) 与机器人操作系统及嵌入式平台 (专用平台/AI)}：正逆运动学控制、机械臂末端轨迹规划。
    \item \textbf{图形核心渲染管线 (GIT) 与并行分布式计算 (PDC)}：超大规模 GPU 细粒度数据级并行流式计算。
\end{itemize}

贯穿全书的核心横切主题包括：\textbf{效能}、\textbf{伦理}、\textbf{建模}、\textbf{编程}、\textbf{原型设计}、\textbf{可用性} 以及 \textbf{综合评估}。

\section{真实能力任务规范}

\begin{taskbox}[任务 1：多场景数字图像格式甄别与工程选型]
\textbf{【任务情境】}：某工业智能装备企业正在开发一套集产品外观全彩高精摄影、二维矢量卡通吉祥物绘制以及数控超高压高精度水刀金属板料切割轨迹规划于一体的多模态系统。\\
\textbf{【能力陈述】}：针对 1) 真实场景彩色摄影照片、2) 扁平化矢量卡通图形、3) 工业数控超高压水刀切割运动路径这三类性质截然不同的媒体数据，准确选定最优的数字图像存储文件格式，并从压缩机制、几何可缩放性与时空开销角度给出严密论证。\\
\textbf{【能力领域】}：应用与理论 \quad \textbf{【能力单元】}：系统评估\\
\textbf{【核心要求】}：图形交互/基础核心 \quad \textbf{【达标等级】}：\textbf{理解与解释}
\end{taskbox}

\begin{taskbox}[任务 2：现代图形渲染管线技术白皮书撰写]
\textbf{【任务情境】}：某新兴游戏引擎研发团队需要为新入职工程师编写一份关于现代 GPU 渲染底层架构的官方技术白皮书。\\
\textbf{【能力陈述】}：以权威、严密的工程语言，系统详尽地阐述从应用程序调用图形 API（如 OpenGL / Vulkan / DirectX）发出绘制调用开始，直至图元最终在物理帧缓冲区完成光栅化与色彩混合输出的完整渲染管线各物理与逻辑阶段。\\
\textbf{【能力领域】}：应用实践 \quad \textbf{【能力单元】}：技术沟通与文档规范\\
\textbf{【核心要求】}：图形交互/基本概念、基础渲染技术 \quad \textbf{【达标等级】}：\textbf{理解与解释}
\end{taskbox}

\begin{taskbox}[任务 3：遵循物理与艺术法则的真实弹跳球动画制作]
\textbf{【任务情境】}：动画工作室正在制作一部高品质三维动画短片，需要构建一个从空中自由跌落并与刚性地面多次反弹碰撞的充气橡胶皮球动态资产。\\
\textbf{【能力陈述】}：综合运用重力加速度、弹性恢复系数、以及经典动画十二原则中的“挤压与拉伸”、“预备动作”、“慢入与慢出（加减速缓冲）”与“时间节拍”，设计并编程或装配渲染出一个物理可信且极具艺术张力的逼真弹跳球动画序列。\\
\textbf{【能力领域】}：工程应用 \quad \textbf{【能力单元】}：需求分析/设计/开发/测试全流程\\
\textbf{【核心要求】}：图形交互/基本概念、基础渲染、计算机动画、高级着色 \quad \textbf{【达标等级】}：\textbf{应用、评估与开发}
\end{taskbox}

\begin{taskbox}[任务 4：基于反应-扩散模型的有机生物条纹生成系统]
\textbf{【任务情境】}：某数字艺术生成实验室需要程序化模拟热带斑马鱼或猎豹皮肤表面的自然有机条纹图案，并要求支持艺术家现场实时操纵化学扩散参数。\\
\textbf{【能力陈述】}：使用自主选定的现代编程语言，编程实现经典的图灵反应-扩散方程（Turing Reaction-Diffusion Model，如 Gray-Scott 偏微分系统），并在屏幕上实时渲染生成的有机条纹；设计并实现一个响应式图形用户界面，使系统的化学物质扩散率、反应速率等参数能够在毫秒级实时联动调节。\\
\textbf{【能力领域】}：应用与开发 \quad \textbf{【能力单元】}：需求分析/架构设计/开发/测试/部署\\
\textbf{【核心要求】}：图形交互/基础渲染、高级着色、物理仿真、交互技术；人机交互/系统设计、用户理解；软件开发基础/开发方法 \quad \textbf{【达标等级】}：\textbf{应用、评估与开发}
\end{taskbox}

\begin{taskbox}[任务 5：高交互可定制的粒子系统数字烟花仿真引擎]
\textbf{【任务情境】}：某虚拟城市文旅庆典平台需要上线一个供全网游客在线自主定制并燃放的 Web 虚拟烟花庆典互动模块。\\
\textbf{【能力陈述】}：编程实现一个包含重力加速度、空气阻力衰减、爆炸冲量分形扩散以及色彩渐变消隐的物理粒子系统烟花发生器；构建直观的人机交互界面，使用户能够实时定制粒子发射速率、色彩调色板、爆炸碎片半径与重力加速度等多维参数。\\
\textbf{【能力领域】}：工程应用 \quad \textbf{【能力单元】}：全周期软件开发与交付\\
\textbf{【核心要求】}：图形交互/基础渲染、物理仿真、交互技术；人机交互/系统设计；软件开发基础/开发方法 \quad \textbf{【达标等级】}：\textbf{应用、评估与开发}
\end{taskbox}

\begin{taskbox}[任务 6：基于高斯卷积核的人类屈光不正视力障碍仿真滤镜]
\textbf{【任务情境】}：眼科医学研究中心与无障碍设计团队需要一套能够向视力健全者直观仿真不同程度近视患者真实视错觉感知的视觉渲染系统。\\
\textbf{【能力陈述】}：设计并实现基于不同核尺寸与标准差 $\sigma$ 的二维离散高斯模糊卷积滤波算法，建立屈光度（近视度数）与卷积核弥散斑之间的精确数学对应关系，对输入的高清物理场景图像进行实时退化模糊仿真。\\
\textbf{【能力领域】}：工程应用 \quad \textbf{【能力单元】}：需求分析/算法设计/工程测试/用户验收\\
\textbf{【核心要求】}：图形交互/基本概念、基础渲染、数字图像处理；软件开发基础/开发方法 \quad \textbf{【达标等级】}：\textbf{应用、评估与开发}
\end{taskbox}

\begin{taskbox}[任务 7：低对比度退化图像的自适应色彩直方图归一化增强]
\textbf{【任务情境】}：恶劣低照度雾霾天气下安防监控摄像头采集的数字图像对比度极度匮乏，关键目标特征被严重掩蔽，亟待自动化提升图像视觉动态范围。\\
\textbf{【能力陈述】}：针对三通道 RGB 彩色图像，编程实现多通道色彩直方图统计、动态范围自适应线性拉伸以及直方图均衡化处理，在消除曝光不足的同时防止过饱和与高频色彩失真。\\
\textbf{【能力领域】}：工程应用 \quad \textbf{【能力单元】}：算法开发与性能验证\\
\textbf{【核心要求】}：图形交互/基本概念、基础渲染、数字图像处理 \quad \textbf{【达标等级】}：\textbf{应用、评估与开发}
\end{taskbox}

\begin{taskbox}[任务 8：跨移动双平台的多人井字棋交互应用工程实现]
\textbf{【任务情境】}：针对 6 岁及以上低龄儿童群体，开发一款界面友好、反应敏捷、兼顾平板与手机触控交互的跨平台益智井字棋游戏。\\
\textbf{【能力陈述】}：面向 iOS 与 Android 移动操作系统平台，完整设计并实现一套具备状态判定、胜负仲裁、悔棋及平局处理逻辑的井字棋移动应用；界面需针对低龄儿童认知特征进行图标与动效优化，支持低时延触控手势响应。\\
\textbf{【能力领域】}：移动软件工程 \quad \textbf{【能力单元】}：需求分析/交互设计/跨平台开发/测试交付\\
\textbf{【核心要求】}：图形交互/基本概念、基础渲染、交互技术；人机交互/系统设计、理解用户；专用平台/移动开发基础、移动平台架构；软件开发基础/开发方法 \quad \textbf{【达标等级】}：\textbf{应用、评估与开发}
\end{taskbox}

\begin{taskbox}[任务 9：区域历史气象气温时空动态演化可视化系统]
\textbf{【任务情境】}：国家气象科学数据中心需要将近半个世纪以来某特大城市群的地表温度时序栅格数据转换为直观生动的动态演化可视化成果，用于联合国气候大会公众科普。\\
\textbf{【能力陈述】}：导入多年连续气温时空观测数据集，设计并实现一套高质量的三维/二维温度时序动态演化动画系统；采用感知均匀的伪彩色色阶编码，叠加等温线动态拓扑形变，生动展示区域气温长期冷暖波动特征。\\
\textbf{【能力领域】}：应用与科学计算 \quad \textbf{【能力单元】}：全流程系统实现与科普可用性评估\\
\textbf{【核心要求】}：图形交互/基本概念、基础渲染、数据可视化；人机交互/系统设计、用户理解 \quad \textbf{【达标等级】}：\textbf{应用、评估与开发}
\end{taskbox}

\begin{taskbox}[任务 10：双区域 COVID-19 传染病时序扩散对比可视化]
\textbf{【任务情境】}：全球流行病预防控制机构急需对两个地理特征相似但防疫政策迥异的代表性行政区域的 COVID-19 历史感染扩散时序数据进行严密可视化对比。\\
\textbf{【能力陈述】}：基于流行病学每日新增感染与累计确诊数据集，设计并构建交互式双联对比动画系统；合理运用双 Y 轴折线、动态热力密度图与时间轴无缝联动刷选，使公共卫生决策者能一目了然洞察两地病毒传播峰值与平台期的动态差异。\\
\textbf{【能力领域】}：应用实践 \quad \textbf{【能力单元】}：需求分析/可视化设计/测试/决策评估\\
\textbf{【核心要求】}：图形交互/基本概念、基础渲染、数据可视化；人机交互/交互设计、受众认知 \quad \textbf{【达标等级】}：\textbf{应用、评估与开发}
\end{taskbox}

\begin{taskbox}[任务 11：弱光自适应智能光感微控制器控制终端系统]
\textbf{【任务情境】}：某智慧农业与仓储物流企业需要部署一批无需人工干预、可在暗光或夜间自动触发照明补偿的超低功耗智能终端节点。\\
\textbf{【能力陈述】}：使用工业标准微控制器（如 Arduino、ESP32 或 STM32）与外围电子元器件，独立搭建包含光敏电阻分压传感回路的硬件电路；编写嵌入式固件程序，实现利用 ADC 定期采样环境照度，并在光强跌破设定阈值时通过驱动外部三极管/MOSFET 稳定点亮高功率发光二极管（LED）。\\
\textbf{【能力领域】}：嵌入式与实体计算 \quad \textbf{【能力单元】}：硬件设计/电路焊接/固件编写/联调部署\\
\textbf{【核心要求】}：图形交互/基本概念、实体与物理计算；专用平台/嵌入式平台；人机交互/硬件设计；社会伦理/物联网隐私 \quad \textbf{【达标等级】}：\textbf{解释、应用、评估与开发}
\end{taskbox}

\section{课程设计与开课规划建议}

联合工作组针对不同专业定位的高校，推荐五套典型的专业课程组织与开课规划方案：

\subsection{方案 1：交互式计算机图形学}
\textbf{【总学时】}：44 学时。
\begin{itemize}
    \item \textbf{涵盖模块}：图形交互/基础渲染技术（40 学时）+ 社会伦理/工程伦理（4 学时）。
    \item \textbf{先修课程要求}：计算机编程基础 II (CS2)、线性代数（熟练掌握仿射变换）、三角学。
    \item \textbf{课程能力目标}：修完本课程的学生应深刻理解并能熟练调用现代图形 API 独立构建计算机图形渲染系统；精准掌握几何三维模型坐标变换、虚拟观察摄影机平截头体定位以及点光源/平行光源的光照计算。
\end{itemize}

\subsection{方案 2：实体与物理计算}
\textbf{【总学时】}：44 学时。
\begin{itemize}
    \item \textbf{涵盖模块}：图形交互/实体与物理计算（30 学时）+ 专用平台/嵌入式平台开发（10 学时）+ 社会伦理/隐私保护与无障碍包容性（4 学时）。
    \item \textbf{先修课程要求}：计算机编程基础 I (CS1)。
    \item \textbf{课程能力目标}：修完本课程的学生应能独立设计并装配基础电子电路，熟练编写微控制器固件；透彻掌握电路极性、欧姆定律、信号模数转换及电子仪器安全操作规程。
\end{itemize}

\subsection{方案 3：数字图像处理}
\textbf{【总学时】}：44 学时。
\begin{itemize}
    \item \textbf{涵盖模块}：图形交互/数字图像处理（30 学时）+ 图形交互/基础渲染技术（10 学时）+ 社会伦理/多元包容、数据隐私与知识产权法（4 学时）。
    \item \textbf{先修课程要求}：计算机编程基础 II (CS2)、线性代数。
    \item \textbf{课程能力目标}：修完本课程的学生应熟练掌握数字图像的数字化采集、空域与频域滤波、色彩增强、数学形态学变换、特征提取、有损/无损压缩编码与磁盘持久化全流程。
\end{itemize}

\subsection{方案 4：数据科学与信息可视化}
\textbf{【总学时】}：47 学时。
\begin{itemize}
    \item \textbf{涵盖模块}：图形交互/数据可视化（30 学时）+ 图形交互/基础渲染技术（10 学时）+ 人机交互/理解用户认知（3 学时）+ 社会伦理/数据伦理、算法偏见与隐私规约（4 学时）。
    \item \textbf{先修课程要求}：计算机编程基础 II (CS2)、线性代数。
    \item \textbf{课程能力目标}：修完本课程的学生应能独立清洗并甄别真实科学或社会学数据集，设计并实现一套符合认知心理学规律、严谨防范误导性修辞、兼具高可用性的交互式可视化软件系统。
\end{itemize}

\subsection{方案 5：计算机物理仿真系统}
\textbf{【总学时】}：44 学时。
\begin{itemize}
    \item \textbf{涵盖模块}：图形交互/计算机物理仿真（30 学时）+ 图形交互/基础渲染技术（10 学时）+ 社会伦理/工程伦理（4 学时）。
    \item \textbf{先修课程要求}：计算机编程基础 II (CS2)、线性代数、基础微积分。
    \item \textbf{课程能力目标}：修完本课程的学生应掌握时间连续物理方程向离散数值积分转化的方法论，能够独立编程构建支持用户实时交互、兼顾物理真实感与艺术可控性的物理与非物理动力学仿真引擎。
\end{itemize}

\vspace{10pt}
\noindent\textbf{【课程规范编写委员会】}
\begin{itemize}
    \item \textbf{主席}：Susan Reiser（北卡罗来纳大学阿什维尔分校，UNC Asheville，美国）
    \item \textbf{核心委员}：
    Erik Brunvand（犹他大学，University of Utah，美国）、
    Kel Elkins（NASA 戈达德太空飞行中心科学可视化实验室，NASA/GSFC，美国）、
    Jeff Lait（SideFX 软件公司，Toronto，加拿大）、
    Amruth Kumar（拉马波学院，Ramapo College，美国）、
    Paul Mihail（瓦尔多斯塔州立大学，Valdosta State University，美国）、
    Tabitha Peck（戴维森学院，Davidson College，美国）、
    Ken Schmidt（美国国家海洋和大气管理局 NCEI，NOAA NCEI，美国）、
    Dave Shreiner（Unity Technologies，旧金山，美国）。
    \item \textbf{主要贡献专家}：
    Greg Shirah（NASA/GSFC 科学可视化实验室）、
    AJ Christensen（NASA/GSFC 科学可视化实验室 / SSAI）、
    Barbara Mones（华盛顿大学，University of Washington）、
    Beatriz Sousa Santos（阿威罗大学，University of Aveiro，葡萄牙）、
    Ted Kim（耶鲁大学，Yale University）、
    Ginger Alford（南卫理公会大学，Southern Methodist University）。
\end{itemize}

\clearpage
\onecolumn
\section{附录：核心主题与技能要求等级对照表}

\noindent\begin{tabularx}{\textwidth}{p{0.15\textwidth} Y p{0.12\textwidth} p{0.10\textwidth} p{0.08\textwidth}}
\toprule
\textbf{知识单元} & \textbf{核心知识点与具体主题} & \textbf{认知等级} & \textbf{分类} & \textbf{学时} \\
\midrule
\multirow{4}{*}{\makecell[l]{图形交互/基本概念}} 
& 娱乐、工商业与科学应用领域的现代图形学落地生态 & 解释 & 计算机核心 & 1 \\
& 人类视觉生理系统、三刺激接收、视差与视觉无障碍 & 解释 & 计算机核心 & 1 \\
& 模拟数据数字化、采样量化、光栅图元 vs 矢量数学模型 & 解释 & 计算机核心 & 1 \\
& 标准媒体格式、加色/减色色彩空间、动画离散时序与伦理 & 解释 & 计算机核心 & 1 \\
\midrule
\multirow{3}{*}{\makecell[l]{图形交互/基础渲染}} 
& 现代渲染管线、仿射空间变换、物理光辐射与数值积分、纹理映射、MIP-mapping、走样与深度消隐 & 解释 & 领域核心 & 10 \\
& 正向与反向渲染范式、多边形网格拓扑、辐射度量、光线追踪、渲染方程、三角形光栅化、空间树结构 & 解释 & 领域核心 & 5 \\
& 基于现代可编程着色器标准 API（如 OpenGL）生成交互图像 & 实现 & 领域核心 & 3 \\
\midrule
\multirow{4}{*}{\makecell[l]{图形交互/数据可视化}} 
& 二维/三维标量场色彩映射、等值面与时变动态数据可视化 & 解释/实现 & 领域核心 & 3 \\
& 多元数据高维降维技术、平行坐标轴、流场与关系拓扑图 & 解释/实现 & 领域核心 & 2 \\
& 驱动视觉符号抽象的认知心理学底座与色彩学原理 & 解释 & 领域核心 & 1 \\
& 可视化设计认知偏见与伦理责任（防范视觉欺骗修辞） & 评估 & 领域核心 & 1 \\
\midrule
\multirow{3}{*}{\makecell[l]{图形交互/几何建模}} 
& 表面模型、细分多边形化、Delaunay 剖分、Marching Cubes 算法、参数多项式曲线曲面（Bézier/NURBS） & 解释/使用 & 领域核心 & 2 \\
& 体素与体积几何模型、SDF 距离场、稀疏体素网格与 CSG 构造实体几何集合布尔运算 & 解释/使用 & 领域核心 & 2 \\
& 过程式几何表征、分形拓扑、L-系统文法派生与生成式三维建模 & 解释/使用 & 领域核心 & 2 \\
\midrule
\makecell[l]{图形交互/高级着色} 
& 渲染方程蒙特卡洛高阶求解、时间动态模糊、镜头景深、级联阴影、遮挡剔除、面光源、Hi-Z、非真实感 NPR & 解释/使用 & 领域核心 & 6 \\
\midrule
\multirow{3}{*}{\makecell[l]{图形交互/计算机动画}} 
& 迪士尼经典动画十二原则（挤压拉伸、时间节奏等）、关键帧时间轴编辑与多维曲线插值（Lerp/Slerp/样条） & 解释/使用 & 领域核心 & 2 \\
& 骨骼装配体系、正向运动学（FK）与逆向运动学（IK）求解 & 解释/使用 & 领域核心 & 2 \\
& 空间几何变换复合链条、旋转多重表征、虚拟摄影机注视动画 & 解释/实现 & 领域核心 & 2 \\
\midrule
\multirow{3}{*}{\makecell[l]{图形交互/物理仿真}} 
& 粒子动力学系统、牛顿力学离散数值积分法（辛欧拉/Verlet） & 解释/实现 & 领域核心 & 2 \\
& 碰撞检测与动态物理响应算法（SDF 求交、分离轴定理） & 解释/实现 & 领域核心 & 2 \\
& 基于欧拉网格的流体动力学、半拉格朗日平流与泊松压力投影 & 解释/实现 & 领域核心 & 2 \\
\midrule
\multirow{3}{*}{\makecell[l]{图形交互/沉浸式技术}} 
& 扩展现实谱系定义辨析（VR/AR/MR）、医学与国防仿真应用 & 解释 & 领域核心 & 1 \\
& 双目立体显示光学、6-DoF 空间位姿追踪、动光延迟抑制 & 解释/实现 & 领域核心 & 3 \\
& 沉浸人因安全（晕动症规避）、无障碍通用设计与空间隐私伦理 & 解释/评估 & 领域核心 & 2 \\
\midrule
\multirow{2}{*}{\makecell[l]{图形交互/交互技术}} 
& 事件驱动编程范式、WIMP 单通道图形用户界面、力反馈外骨骼 & 解释/实现 & 领域核心 & 4 \\
& 图形用户界面无障碍等价通道设计与多模态脑机前沿评估 & 解释/评估 & 领域核心 & 2 \\
\midrule
\multirow{3}{*}{\makecell[l]{图形交互/数字图像处理}} 
& 空间离散卷积滤波核（高斯平滑、拉普拉斯边缘锐化） & 解释/实现 & 领域核心 & 2 \\
& 卷积神经网络与视觉 Transformer 架构（与 AI 领域交叉） & 解释/实现 & 领域核心 & 2 \\
& 色彩直方图均衡化、正交频域变换（傅里叶/离散余弦变换） & 解释/实现 & 领域核心 & 2 \\
\midrule
\multirow{4}{*}{\makecell[l]{图形交互/实体与物理计算}} 
& 与真实物理世界通信交互、传感器 ADC 采集与执行器 PWM 驱动 & 解释/实现 & 领域核心 & 1 \\
& 微控制器事件驱动固件架构、硬件定时器与外部中断响应 & 解释/实现 & 领域核心 & 1 \\
& 连接物理实体、CAD/CAM 数字制造（3D 打印/数控机床/吸塑成型） & 解释/评估/实现 & 领域核心 & 3 \\
& 物联网体系架构、轻量级网络通信协议与嵌入式物理安全 & 解释 & 领域核心 & 1 \\
\bottomrule
\end{tabularx}
\clearpage
\twocolumn